\songtitle{A Borrowed Time}

\tag*{2026}\tag{folk-rock}\tag{autobiographical}\label{sec:a-borrowed-time}
\begin{notes}
Lyrics by Suno based on a Claude-produced excerpt of \emph{A Life in World Events}, the memoir of Carlos Thompson -- Braunschweig, July 1989, sharing a hostel with teenagers from Kuwait during a contest; a year later, at a school in Stockholm with Iraqi classmates, the news of the invasion of Kuwait: two sets of friends inside one date. Retrieved from Carlos Th's Suno page (V6-MINI, 12 November 2026).
\end{notes}
\begin{style}
Mid-tempo folk-rock ballad, intimate and story-led. Warm acoustic guitar establishes a steady pulse; restrained electric guitar, rounded bass, brushed snare, and subtle organ gradually widen the arrangement. Male baritone, close and conversational in the verses, rising into a memorable but understated chorus. Begin with a brief spoken news-anchor-style call, then let the song unfold chronologically; pull back for the Stockholm classroom verse, build gently into the televised-war announcement, and finish with a quiet reflective coda. Organic, clear production with natural room detail and no bombastic arena finish.
\end{style}
\begin{lyrics}
\songpart{Spoken Intro}
This is a news bulletin. In August 1990, Iraq has invaded Kuwait.

\songpart{Verse 1}
Braunschweig, July ’89,\\
A few miles from the border line,\\
We waited for the contest to begin,\\
With borrowed keys and numbered rooms.\\
Colombia was down the hall,\\
Kuwait came through the breakfast room,\\
We traded names, we passed the time,\\
Like teenagers do when days are wide.

\songpart{Chorus}
We were only there together,\\
For a little while, a little while,\\
Sharing corridors and morning bread,\\
Before the world would name the miles.\\
We had no map for what was coming,\\
No warning in the summer air;\\
We were friends beneath one hostel roof,\\
And none of us knew what waited there.

\songpart{Verse 2}
By November, lines would change,\\
The old border would not mean the same,\\
But in that July, nothing told us so,\\
We packed our cases, watched the hours go.\\
And nothing spoke of what would soon\\
Break over Kuwait beneath another moon;\\
We waved goodbye, a year still far away,\\
With ordinary plans for ordinary days.

\songpart{Chorus}
We were only there together,\\
For a little while, a little while,\\
Sharing corridors and morning bread,\\
Before the world would name the miles.\\
We had no map for what was coming,\\
No warning in the summer air;\\
We were friends beneath one hostel roof,\\
And none of us knew what waited there.

\songpart{Verse 3}
A year later, Stockholm, after Beijing,\\
A levelling school, new faces gathering:\\
Poland, Ethiopia, Eritrea,\\
Zaire, Angola, Czechoslovakia;\\
Iraq, Vietnam, Iran in the room,\\
Bolivia, Mexico, Spain, too.\\
Then August came across the classroom floor:\\
Iraq had crossed into Kuwait’s door.

\songpart{Bridge}
I knew the Iraqi kids beside me,\\
Knew the Kuwaiti friends from that old stay.\\
Two groups of teenagers, once sharing space,\\
Now pulled to different sides by history’s pace.\\
The news was not a far-off country’s name;\\
It sat among the desks, and felt the same.

\songpart{Verse 4}
Later came the ultimatum, later still,\\
The anchor said the screens would show it live:\\
The first war fully televised, they said,\\
As CNN carried it into every home.\\
But first there was the classroom, still and close,\\
The faces I had known, the fact of it—\\
No broadcast could explain that sudden weight:\\
Two sets of friends inside one date.

\songpart{Final Chorus}
We were only there together,\\
For a little while, a little while,\\
Sharing corridors and morning bread,\\
Before the world would name the miles.\\
We had no map for what was coming,\\
No warning in the summer air;\\
We were friends beneath one hostel roof,\\
And none of us knew what waited there.

\songpart{Outro}
Braunschweig, July ’89,\\
A breakfast room, a borrowed time.\\
We said goodbye and went our ways—\\
The world changed just beyond those days.
\end{lyrics}

